\documentclass[11pt]{article}
\usepackage[T1]{fontenc}
\usepackage[utf8]{inputenc}
\usepackage{lmodern}
\usepackage[a4paper,margin=27mm]{geometry}
\usepackage{amsmath,amssymb,amsthm,booktabs,array}
\usepackage[numbers,sort&compress]{natbib}
\usepackage{microtype}
\usepackage[colorlinks=true,linkcolor=blue,citecolor=blue,urlcolor=blue]{hyperref}
\newcommand{\ket}[1]{\lvert #1\rangle}
\newcommand{\bra}[1]{\langle #1\rvert}
\newcommand{\Tr}{\operatorname{Tr}}
\newcommand{\id}{\operatorname{id}}
\newcommand{\N}{\mathcal N}
\newcommand{\E}{\mathcal E}
\newcommand{\K}{\mathcal K}
\newtheorem{theorem}{Theorem}
\title{An exact anti-degradability threshold for the\newline three-level bosonic loss--dephasing channel}
\author{Xiaojian Zeng\\[3pt]
\small PBC School of Finance, Tsinghua University, 100190\\
\small \href{mailto:xiaojianzeng@foxmail.com}{xiaojianzeng@foxmail.com}}
\date{}
\begin{document}
\maketitle
\begin{abstract}
We determine the exact anti-degradability threshold of the bosonic
loss--dephasing channel restricted to its first three Fock levels, at Gaussian
phase variance $\ln 2$. The threshold is an isolated root of a sextic
polynomial, approximately $0.664544268033$. A rank-four symmetric Choi
extension establishes feasibility at the endpoint. Explicit separating
operators, lifted from the population support to the full extension space,
exclude every larger transmissivity. Exact loss composition covers the remaining
parameter intervals. This certifies a previously numerical boundary and
sharpens a necessary condition for anti-degradability of the unrestricted
bosonic channel.
\end{abstract}
\noindent\textbf{Keywords:} bosonic quantum channel; photon loss; Gaussian dephasing;
anti-degradability; symmetric extension

\section{Introduction}
Photon loss and random phase noise act on different parts of a bosonic state:
loss changes photon-number populations, while number dephasing suppresses
coherences. Their combined channel can be anti-degradable even in a parameter
region where pure loss alone is not. An anti-degradable channel has zero
unassisted quantum capacity, because its environment can simulate the receiver.
Determining the boundary of this property therefore provides a useful,
although incomplete, communication criterion.

Mele, Salek, Giovannetti and Lami studied the bosonic loss--dephasing channel
and its finite Fock restrictions \cite{Mele2024}. Their two-level threshold is
$1/(1+e^{-\gamma})$. They also established equality with this bound for the
three-level restriction when $e^{-\gamma}\leq\sqrt2-1$; outside that regime
the three-level boundary was evaluated numerically. We fix $\gamma=\ln2$,
where that analytic hypothesis does not apply. Our contribution is the exact
endpoint and explicit primal and full-space dual certificates for this
previously numerical point. We do not claim that the qualitative departure
from the two-level boundary was previously unknown.

The argument exploits zeros in the channel's population support and its
number-phase covariance. Swap invariance alone does not justify restricting
an extension to the symmetric output subspace in arbitrary qutrit problems
\cite{Myhr2009}. Here that restriction follows from a particular
marginal-preserving transfer of antisymmetric blocks. A positive operator
constructed on the resulting support is then lifted to the entire
$3\times3\times3$ extension space. This distinction is essential for an
unrestricted infeasibility certificate.

\section{Channel, conventions and result}
Let $\mathcal H_3=\operatorname{span}\{\ket0,\ket1,\ket2\}$. The exact
restriction of vacuum pure loss followed by centered Gaussian number
dephasing is
\begin{equation}\label{eq:channel}
\begin{split}
\N_\lambda(\ket m\bra n)
={}&2^{-(m-n)^2/2}\sum_{\ell=0}^{\min(m,n)}
\sqrt{\binom m\ell\binom n\ell}\,(1-\lambda)^\ell
\lambda^{(m+n)/2-\ell}\ket{m-\ell}\bra{n-\ell},
\end{split}
\end{equation}
for $m,n\in\{0,1,2\}$ and transmissivity $0\leq\lambda\leq1$.
The dephasing convention is $e^{-\gamma(m-n)^2/2}$ with phase variance
$\gamma=\ln2$, as in Lemma 25 of Ref.~\cite{Mele2024}. The parameters,
phases in radians, photon numbers and all matrix entries are dimensionless.
This is an invariant finite-dimensional restriction: loss only removes
photons. There is no output projection, postselection or renormalization,
and no perturbative or numerical truncation of the channel.

In the input-first tensor convention, define the normalized Choi state
\begin{equation}\label{eq:choi}
J_\lambda=\frac13\sum_{m,n=0}^2\ket m\bra n\otimes
\N_\lambda(\ket m\bra n),\qquad \Tr_BJ_\lambda=I_A/3.
\end{equation}
Writing $\ket{ab}=\ket a_A\otimes\ket b_B$, its nonzero entries are
\begin{equation}\label{eq:choi-entries}
\begin{split}
3J_\lambda={}&\ket{00}\bra{00}+(1-\lambda)\ket{10}\bra{10}
+\lambda\ket{11}\bra{11}\phantom{\Big|}\\
&+(1-\lambda)^2\ket{20}\bra{20}
+2\lambda(1-\lambda)\ket{21}\bra{21}
+\lambda^2\ket{22}\bra{22}\\
&+\left[\sqrt{\lambda/2}\ket{00}\bra{11}
+\frac\lambda4\ket{00}\bra{22}
+\sqrt\lambda(1-\lambda)\ket{10}\bra{21}
+\frac{\lambda^{3/2}}{\sqrt2}\ket{11}\bra{22}
+\mathrm{h.c.}\right].
\end{split}
\end{equation}
No partial transpose is included. A channel is anti-degradable precisely
when its Choi state has an extension $X$ satisfying
\begin{equation}\label{eq:extension}
X\succeq0,\qquad SX S=X,\qquad \Tr_{B_2}X=J_\lambda,
\end{equation}
where $S$ exchanges $B_1$ and $B_2$ \cite{Mele2024,Myhr2009}.
Swap symmetry gives the other marginal. For completeness, an anti-degrading
map applied to a Stinespring environment gives a channel with two identical
output marginals; swap averaging gives \eqref{eq:extension}. Conversely,
the extension defines such a two-output channel. Uniqueness of Stinespring
dilation expresses its second output as a channel from the environment
of the first, providing an anti-degrading map.

\begin{theorem}\label{thm:threshold}
Let $r$ be the unique root in
\begin{equation}\label{eq:root-interval}
I_r=\left[\frac{66454426803}{10^{11}},
                \frac{66454426804}{10^{11}}\right]
\end{equation}
of
\begin{equation}\label{eq:polynomial}
P(x)=601x^6-714x^5+113x^4-1040x^3+2080x^2-1280x+256.
\end{equation}
Then $\N_\lambda$ is anti-degradable if and only if $\lambda\leq r$.
At $r$ there is a rank-four normalized swap-invariant extension.
For every $\lambda>r$ there is an explicit Hermitian $H_\lambda$ such that
\begin{equation}\label{eq:dual-condition}
W_{H_\lambda}:=\frac{H_\lambda\otimes I_3+
S(H_\lambda\otimes I_3)S}{2}\succeq0,\qquad
\Tr(H_\lambda J_\lambda)<0
\end{equation}
on the full 27-dimensional space.
\end{theorem}

An exact narrower enclosure is
\begin{equation}\label{eq:refined}
\frac{66454426803317}{10^{14}}<r<
\frac{66454426803318}{10^{14}}.
\end{equation}
Thus $r\simeq0.664544268033$, slightly below the two-level bound $2/3$.
The decimal is descriptive; the proof uses the algebraic definition and
rational enclosures.

\section{Population support and the two swap sectors}\label{sec:reduction}
Set $Y=3X$, so that $\Tr_{B_2}Y=3J_\lambda$ and $\Tr Y=3$.
All marginal populations with $b>a$ vanish. For a positive $Y$, the
corresponding diagonal summands $Y_{abc,abc}$ are nonnegative and hence
individually zero. A zero diagonal of a positive matrix forces its whole
row and column to vanish, by $|Y_{ij}|^2\leq Y_{ii}Y_{jj}$.
The second marginal imposes $c\leq a$ as well. Every extension therefore
has support in
\begin{equation}\label{eq:support}
\K=\operatorname{span}\{\ket{abc}:b,c\leq a\},\qquad
\dim\K=1+4+9=14.
\end{equation}

The Choi state is invariant under $e^{i\theta(-N_A+N_B)}$.
To average an extension, use the representation
\begin{equation}
U_\theta=e^{i\theta(-N_A+N_{B_1}+N_{B_2})}.
\end{equation}
This preserves positivity and each marginal and gives blocks of charge
$q=b+c-a$. The averaging commutes with the swap. For $b<c$ define
\begin{equation}
\ket{a;s_{bc}}=\frac{\ket{abc}+\ket{acb}}{\sqrt2},\qquad
\ket{a;a_{bc}}=\frac{\ket{abc}-\ket{acb}}{\sqrt2};
\quad\ket{a;s_{bb}}=\ket{abb}.
\end{equation}
The symmetric sector has charge dimensions $(1,2,4,2,1)$, and the
antisymmetric sector has dimensions $(1,2,1)$ at $q=-1,0,1$.

For this support, every antisymmetric block can be transferred isometrically
to the corresponding symmetric kets without changing either marginal.
The singleton kets are $\ket{2;a_{01}}$ and $\ket{2;a_{12}}$.
Their diagonal marginal images coincide with those of the symmetric kets.
The only nonsingleton block has $\ket{1;a_{01}}$ and
$\ket{2;a_{02}}$. Its two diagonal units have the same property.
In its off-diagonal unit the only common traced number is zero;
the two antisymmetric signs agree, so the partial trace equals that
of the symmetric unit. The same holds for its adjoint and for the
other partial trace. These six ordered matrix units exhaust all blocks.
Adding the transferred positive matrices to the symmetric blocks preserves
positivity. This proves symmetric-support equivalence for the present
channel; no such equivalence for general qutrit channels is used.

Since $J_\lambda$ is real, averaging with complex conjugation permits
real symmetric blocks. Their ordered bases are in Table~\ref{tab:bases}.
Let $B_q$ denote the isometry whose columns are the listed kets.
\begin{table}[htbp]
\centering
\caption{Ordered symmetric block bases in the necessary support.
The scalar blocks are denoted $\zeta$ and $w$.}\label{tab:bases}
\begin{tabular}{ccl}
\toprule
Charge & Block & Ordered basis\\
\midrule
$-2$ & $\zeta$ & $\ket{2;s_{00}}$\\
$-1$ & $R$ & $\ket{1;s_{00}},\ \ket{2;s_{01}}$\\
$0$ & $T$ & $\ket{0;s_{00}},\ \ket{1;s_{01}},\ \ket{2;s_{02}},\ \ket{2;s_{11}}$\\
$1$ & $V$ & $\ket{1;s_{11}},\ \ket{2;s_{12}}$\\
$2$ & $w$ & $\ket{2;s_{22}}$\\
\bottomrule
\end{tabular}
\end{table}
Direct partial tracing of these kets gives the ten independent equations
\begin{equation}\label{eq:marginals}
\begin{aligned}
T_{00}&=1,& T_{01}&=\sqrt\lambda,& T_{02}&=\lambda/(2\sqrt2),\\
R_{00}+T_{11}/2&=1-\lambda,& T_{11}/2+V_{00}&=\lambda,\\
R_{01}+T_{13}&=\sqrt{2\lambda}(1-\lambda),\\
T_{12}+\sqrt2 V_{01}&=\sqrt2\lambda^{3/2},\\
\zeta+R_{11}/2+T_{22}/2&=(1-\lambda)^2,\\
R_{11}/2+T_{33}+V_{11}/2&=2\lambda(1-\lambda),\\
T_{22}/2+V_{11}/2+w&=\lambda^2.
\end{aligned}
\end{equation}
All other marginal entries vanish. The problem now has 18 real matrix
entries, with ten independent constraints, while retaining equivalence
to \eqref{eq:extension}.

\section{The algebraic endpoint and its extension}\label{sec:primal}
At the specified root $r$, set
\begin{equation}\label{eq:primal-parameters}
t=\sqrt r,\quad d=2-3r,\quad p=\frac{r}{2\sqrt2},\quad
\alpha=2(1-r)^2-\frac{r^2}{8},\quad k=\frac{t}{\sqrt2}.
\end{equation}
All square roots below are the positive real roots. Set $\zeta=0$, $R=0$,
and
\begin{equation}\label{eq:primal}
G=\begin{pmatrix}
1&0\\ t&\sqrt d\\ p&\sqrt\alpha\\ kt&k\sqrt d
\end{pmatrix},\qquad
T=GG^{\mathsf T},\qquad
h=\begin{pmatrix}\sqrt{2r-1}\\\sqrt{2r(1-r)}\end{pmatrix},\qquad
V=hh^{\mathsf T},\quad w=r^2+r-1.
\end{equation}
The positivity of $d,\alpha,2r-1,w$ follows from the rational interval
\eqref{eq:root-interval}; exact bounds are summarized in Appendix~\ref{app:arithmetic}.
Thus these blocks are positive semidefinite. Substitution in
\eqref{eq:marginals} leaves only one residual equation,
\begin{equation}\label{eq:radical}
\sqrt{(2-3r)\bigl(2(1-r)^2-r^2/8\bigr)}
+2\sqrt{r(1-r)(2r-1)}=\frac{3r^{3/2}}{2\sqrt2}.
\end{equation}

To check the branch after squaring, define
\begin{equation}\label{eq:ABC}
\begin{split}
A(x)&=(2-3x)\bigl(16(1-x)^2-x^2\bigr),\\
B(x)&=32x(1-x)(2x-1),\qquad C(x)=9x^3.
\end{split}
\end{equation}
Direct expansion gives
\begin{equation}\label{eq:squaring}
(C-A-B)^2-4AB=4P(x).
\end{equation}
On $I_r$, $A,B,C,C-A-B$ are strictly positive. Consequently $P(r)=0$
implies $C-A-B=2\sqrt{AB}$ and then
$\sqrt A+\sqrt B=\sqrt C$, which is \eqref{eq:radical}.
There is no extraneous negative radical branch. Sturm counts establish
exactly one root in $I_r$ and exactly one in $(3/5,2/3)$.
Applying the same exact count to \eqref{eq:refined} gives its narrower
enclosure.

The endpoint extension is now explicit:
\begin{equation}\label{eq:endpoint-X}
X_r=\frac13\left[B_0GG^{\mathsf T}B_0^{\mathsf T}
+B_1hh^{\mathsf T}B_1^{\mathsf T}
+w\ket{222}\bra{222}\right].
\end{equation}
Equations~\eqref{eq:marginals}--\eqref{eq:radical} give both marginals
$J_r$. Each component is supported on symmetric output kets.
Moreover $\Tr T+\Tr V+w=3$, so $\Tr X_r=1$.
The two columns of $G$ are independent because $d>0$; $h\ne0$ and $w>0$.
The three charge sectors are orthogonal, proving $\operatorname{rank}X_r=4$.

\section{A separating operator on the full population support}\label{sec:face}
The following coefficients are fixed at $r$, even when the tested
transmissivity is varied. Define
\begin{equation}\label{eq:dual-parameters}
u=\sqrt{\alpha/d},\quad
v=\sqrt{\frac{2r(1-r)}{2r-1}},\quad
\eta=\frac{\sqrt2u}{v},\quad \beta=2\eta,\quad c=tu-p,
\end{equation}
and ten real numbers
\begin{equation}\label{eq:y}
\begin{aligned}
y_0&=c^2,& y_1&=-2\sqrt2uc,& y_2&=2\sqrt2c,\\
y_3&=2u^2+\eta(2r-v^2),& y_4&=-2\beta t,& y_5&=\eta v^2,\\
y_6&=-4u,& y_7&=2,&y_8&=\beta,&y_9&=0.
\end{aligned}
\end{equation}
In the lexicographic basis $\ket{00},\ket{01},\ldots,\ket{22}$,
let $H_0$ have diagonal entries $y_0,y_3,y_5,y_7,y_8,y_9$ at
indices $0,3,4,6,7,8$, respectively. Its only nonzero off-diagonal
entries and their equal adjoints are
\begin{equation}\label{eq:H0}
(H_0)_{0,4}=y_1/2,\quad (H_0)_{0,8}=y_2/2,\quad
(H_0)_{3,7}=y_4/2,\quad (H_0)_{4,8}=y_6/2.
\end{equation}
Define $W_0=W_{H_0}$ as in \eqref{eq:dual-condition}.
Its symmetric blocks on $\K$ are
\begin{equation}\label{eq:dual-blocks}
\begin{split}
C_\zeta&=2,\qquad
C_R=\begin{pmatrix}y_3&-\beta k\\-\beta k&1+\eta\end{pmatrix},\\
C_T&=nn^{\mathsf T}+\beta mm^{\mathsf T},\qquad
C_V=\eta jj^{\mathsf T},\qquad C_w=0,\\
n&=(c,-u,1,0)^{\mathsf T},\quad
m=(0,-k,0,1)^{\mathsf T},\quad j=(v,-1)^{\mathsf T}.
\end{split}
\end{equation}
These identities follow by compression in the bases of Table~\ref{tab:bases}.
For example $(C_T)_{12}=y_6/4$ and $(C_V)_{01}=y_6/(2\sqrt2)$:
the factors account for the two real off-diagonal entries in a trace pairing.
Exact interval bounds give $y_3>53$ and $334<\det C_R<335$.
Thus $C_R$ is positive definite and the remaining blocks are manifestly
positive semidefinite. The antisymmetric blocks have identical marginal
images to the corresponding symmetric principal blocks, by
Section~\ref{sec:reduction}. Their trace pairing with $H_0$ is therefore
identical, so they too are positive semidefinite. This proves
$P_\K W_0P_\K\succeq0$ on all 14 support dimensions.

Complementary slackness can be checked without an eigenvalue calculation:
$n^{\mathsf T}G=m^{\mathsf T}G=0$ and $j^{\mathsf T}h=0$.
Consequently $q(r)=0$, where the variable-transmissivity expectation is
\begin{equation}\label{eq:q}
\begin{split}
q(\lambda):=3\Tr(H_0J_\lambda)
={}&y_0+\frac{y_1\sqrt\lambda}{\sqrt2}+\frac{y_2\lambda}{4}
+y_3(1-\lambda)+y_4\sqrt\lambda(1-\lambda)+y_5\lambda\\
&+\frac{y_6\lambda^{3/2}}{\sqrt2}+y_7(1-\lambda)^2
+2y_8\lambda(1-\lambda)+y_9\lambda^2.
\end{split}
\end{equation}
The coefficients remain fixed in this differentiation:
\begin{equation}\label{eq:qprime}
\begin{split}
q'(\lambda)={}&\frac{y_1}{2\sqrt2\sqrt\lambda}+\frac{y_2}{4}-y_3
+\frac{y_4(1-3\lambda)}{2\sqrt\lambda}+y_5
+\frac{3y_6\sqrt\lambda}{2\sqrt2}\\
&-2y_7(1-\lambda)+2y_8(1-2\lambda).
\end{split}
\end{equation}
Rational outward interval evaluation on $[\inf I_r,2/3]$ yields
$-83<q'(\lambda)<-81$, as detailed in Appendix~\ref{app:arithmetic}.
Hence $q(\lambda)<0$ for every $r<\lambda\leq2/3$.
This statement is uniform over the interval, rather than a test of sample points.

\section{Full-space witnesses and the remaining transmissivities}\label{sec:lift}
Positivity of $W_0$ outside $\K$ has not been assumed. Introduce
\begin{equation}\label{eq:penalty}
F=\operatorname{diag}(0,1,1,0,0,1,0,0,0).
\end{equation}
It marks the forbidden output populations $b>a$ and satisfies
$\Tr(FJ_\lambda)=0$ for every $\lambda$. The associated $W_F$
vanishes on $\K$ and is diagonal with entries at least $1/2$ on
$\K^\perp$. The rational coefficient bounds give $|y_i|<60$;
the 27-dimensional row-sum bound then gives
$\lVert W_0\rVert_{\mathrm{op}}<27\cdot60< K_0:=2000$.

For $r<\lambda\leq2/3$ define
\begin{equation}\label{eq:full-H}
\epsilon_\lambda=-q(\lambda)/6>0,\qquad
M_\lambda=2\left(K_0+\frac{K_0^2}{\epsilon_\lambda}\right),\qquad
H_\lambda=H_0+\epsilon_\lambda I_9+M_\lambda F.
\end{equation}
In $\K\oplus\K^\perp$, write $W_{H_\lambda}$ in block form
$\bigl(\begin{smallmatrix}A_*&B_*\\B_*^\dagger&D_*\end{smallmatrix}\bigr)$.
We have $A_*\succeq\epsilon_\lambda I$, $\lVert B_*\rVert\leq K_0$,
and $D_*\succeq(-K_0+\epsilon_\lambda+M_\lambda/2)I$.
Therefore
\begin{equation}\label{eq:schur}
\begin{split}
D_*-B_*^\dagger A_*^{-1}B_*
&\succeq\left(-K_0+\epsilon_\lambda+\frac{M_\lambda}{2}
-\frac{K_0^2}{\epsilon_\lambda}\right)I
=\epsilon_\lambda I>0.
\end{split}
\end{equation}
Thus the full $W_{H_\lambda}$ is positive definite, whereas
\begin{equation}\label{eq:negative}
\Tr(H_\lambda J_\lambda)=q(\lambda)/3+\epsilon_\lambda
=q(\lambda)/6<0.
\end{equation}
For any extension satisfying \eqref{eq:extension},
$\Tr(W_{H_\lambda}X)=\Tr(H_\lambda J_\lambda)$ would be nonnegative,
a contradiction. The penalty grows without bound as $\lambda\downarrow r$,
but is finite for each strictly larger transmissivity.

Let $\E_\mu$ be vacuum pure loss restricted to $\mathcal H_3$.
Its Kraus operators are
\begin{equation}\label{eq:loss-kraus}
L_\ell(\mu)=\sum_{n=\ell}^2\sqrt{\binom n\ell}
(1-\mu)^{\ell/2}\mu^{(n-\ell)/2}\ket{n-\ell}\bra n,
\quad\ell=0,1,2.
\end{equation}
They satisfy $\sum_\ell L_\ell^\dagger L_\ell=I_3$ throughout $[0,1]$.
Loss and number dephasing commute, and the exact loss composition law gives
\begin{equation}\label{eq:composition}
\E_\mu\circ\N_\lambda=\N_{\mu\lambda}
\end{equation}
(Lemma 28 and Eq.~(67) in Ref.~\cite{Mele2024}). It can also be checked
directly on the nine matrix units using \eqref{eq:channel} and
\eqref{eq:loss-kraus}.

For $0\leq\lambda\leq r$, set $\Psi=\E_{\lambda/r}$ and
\begin{equation}\label{eq:below}
X_\lambda=(\id_A\otimes\Psi\otimes\Psi)(X_r).
\end{equation}
Complete positivity, trace preservation and applying the same map to both
outputs preserve positivity, normalization and swap invariance. Its marginals
are $J_\lambda$ by \eqref{eq:composition}.

For $2/3\leq\lambda\leq1$, take $U=2/3$ and the operator $H_U$ from
\eqref{eq:full-H}. With $\Psi=\E_{U/\lambda}$, define
\begin{equation}\label{eq:pullback}
H_\lambda=(\id_A\otimes\Psi^\dagger)(H_U).
\end{equation}
The adjoint is completely positive and unital, giving
\begin{equation}\label{eq:pullback-W}
\begin{split}
W_{H_\lambda}&=(\id_A\otimes\Psi^\dagger\otimes\Psi^\dagger)(W_{H_U})
\succeq0,\\
\Tr(H_\lambda J_\lambda)&=\Tr(H_UJ_U)<0.
\end{split}
\end{equation}
The unital identity is needed here for both tensor terms in $W$.
Together \eqref{eq:endpoint-X}, \eqref{eq:full-H}, \eqref{eq:below} and
\eqref{eq:pullback} establish every assertion of Theorem~\ref{thm:threshold},
including both physical endpoints $\lambda=0,1$.

\section{Physical consequence, comparison and limits}
The exact feasible interval is $[0,r]$ for arbitrary states supported on the
first three Fock levels. In particular the unassisted quantum capacity of
this restriction vanishes there. The two-level threshold $2/3$ is larger
by
\begin{equation}\label{eq:delta}
2/3-r\simeq0.00212239863349.
\end{equation}
By Corollary 35 of Ref.~\cite{Mele2024}, non-anti-degradability of a finite
restriction implies non-anti-degradability of the unrestricted channel.
Thus anti-degradability of the infinite bosonic channel at this phase
variance necessarily requires $\lambda\leq r$, sharpening the two-level
necessary bound. Anti-degradability below $r$ for the unrestricted channel
does not follow from this restriction, and neither does positive quantum
capacity above $r$.

The closest prior calculation is the three-level numerical boundary in
Ref.~\cite{Mele2024}; the advance is exact algebraic and operator
certification at one specified variance. The symmetry-class criterion of
Ranade \cite{Ranade2009} concerns invariance under the full family of
diagonal unitaries. We use only number-phase covariance and do not apply
that criterion to the present Choi matrix. Our explicit construction and
separation argument supply the threshold directly. The literature comparison
supports this narrow contribution; exhaustive priority under every possible
transformed representation is not asserted.

This is an exact finite-restriction result, with rational error control for
the algebraic signs. It is not a general-$\gamma$ formula, a capacity
formula, a two-way communication boundary or an empirical claim.
The support transfer and full-space penalty exhibit how a singular
symmetric-extension endpoint can be certified without treating numerical
solver feasibility as a proof.

\appendix
\section{Exact arithmetic and reproducibility}\label{app:arithmetic}
The auxiliary sign checks use rational endpoints, exact polynomial
arithmetic and outward bounds for square roots. For a nonnegative rational
$a$, with $D=10^{30}$, compute
\begin{equation}
n=\left\lfloor\sqrt{\left\lfloor aD^2\right\rfloor}\right\rfloor.
\end{equation}
Then $n/D\leq\sqrt a<(n+1)/D$. Apply this to the lower and upper
endpoints of a nonnegative interval. Addition and multiplication use
the endpoint extrema; division is performed only when the denominator
interval excludes zero. All such arithmetic uses integers and rational
numbers, with no floating-point tolerance.

Table~\ref{tab:bounds} lists conservative rational enclosures, written
with terminating decimals interpreted as exact rationals. The sign and
operator conditions follow with substantial margins. The coefficient
intervals in \eqref{eq:y} are computed with $r\in I_r$, and the derivative
interval uses $\lambda\in[\inf I_r,2/3]$ independently. This enlarged
interval covers every required value; no cancellation requiring numerical
correlation between the two variables is assumed.
\begin{table}[htbp]
\centering
\caption{Conservative exact rational bounds sufficient for the radical
branch, endpoint positivity and full-space witness construction.}
\label{tab:bounds}
\begin{tabular}{lcc}
\toprule
Quantity & Lower bound & Upper bound\\
\midrule
$A(r)$ & $0.0086$ & $0.0087$\\
$B(r)$ & $2.34$ & $2.35$\\
$C(r)$ & $2.64$ & $2.65$\\
$C(r)-A(r)-B(r)$ & $0.28$ & $0.29$\\
$d$ & $0.0063$ & $0.0064$\\
$\alpha$ & $0.169$ & $0.171$\\
$2r-1$ & $0.329$ & $0.330$\\
$w$ & $0.10$ & $0.12$\\
$y_3$ & $53$ & $54$\\
$\det C_R$ & $334$ & $335$\\
Each $y_i$ & $-60$ & $60$\\
$q'(\lambda)$ on $[\inf I_r,2/3]$ & $-83$ & $-81$\\
\bottomrule
\end{tabular}
\end{table}

For root isolation, take the Sturm chain $S_0=P$, $S_1=P'$,
$S_{j+1}=-\operatorname{rem}(S_{j-1},S_j)$, stopping at the nonzero
constant. Exact sign variations at the endpoints of $I_r$ are respectively
3 and 2, so the root count is one. The same computation gives one root
in $(3/5,2/3)$ and in the interval of \eqref{eq:refined}. These checks,
the polynomial identity \eqref{eq:squaring}, endpoint marginal identities,
dual block identities and the nine loss-composition units are supplied as
reproducible symbolic calculations in Online Resource~1. Numerical
optimization was used only to discover the candidate, not to validate it.

The recorded computational environment is Python 3.12.14, SymPy 1.14.0,
NumPy 2.3.5 and SciPy 1.18.1. The decisive symbolic and rational checks
require SymPy and the Python standard library; NumPy and SciPy support
the optional discovery calculations. The supplementary README separates
these tasks and gives regeneration commands.

\section*{Statements and Declarations}
\paragraph{Data and code availability.}
The channel is specified analytically by \eqref{eq:channel}.
Own symbolic verification code, exact certificate data and reproduction
instructions accompany the manuscript as Online Resource~1; no experimental
dataset is used. The supplementary material does not include third-party
full-text papers.
\paragraph{Use of generative artificial intelligence.}
OpenAI Codex was used in analytic candidate development, calculation and
code preparation, computational checking, literature screening and manuscript
drafting. The exact checks use the symbolic and rational arithmetic methods
described above. This disclosure covers substantive research and text
generation, beyond copy editing.

\bibliographystyle{unsrtnat}
\bibliography{references}
\end{document}
